% !TEX root = task-6-11-2.tex
% Editable task bank. Append new BankTask blocks at the end of the document.
\documentclass[12pt,a4paper]{article}
\usepackage{physics-textbook}
\usepackage{graphicx,needspace}
\usetikzlibrary{decorations.pathmorphing}
\newif\ifShowSolutions
\ifdefined\TutorEdition
  \ShowSolutionstrue
\else
  \ShowSolutionsfalse % Change to \ShowSolutionstrue to display solutions.
\fi
\newcounter{banktask}
\newsavebox{\BankTaskBox}
\newcommand{\BankTaskContent}[3]{%
  \noindent\textbf{\thebanktask.}\quad #2\par
  \ifShowSolutions
    \Needspace{50mm}% Keep the solution heading with its opening text or diagram.
    #3\par
    \ifdefined\TutorEdition
      \ifcsname TutorSource#1\endcsname
        \nopagebreak\smallskip{\small\csname TutorSource#1\endcsname\par}%
      \fi
    \fi
  \fi
}
\newcommand{\BankTask}[3]{%
  \refstepcounter{banktask}\par\medskip
  \sbox{\BankTaskBox}{\begin{minipage}{\linewidth}%
    \BankTaskContent{#1}{#2}{#3}\end{minipage}}%
  \ifdim\dimexpr\ht\BankTaskBox+\dp\BankTaskBox\relax>\dimexpr\textheight-20mm\relax
    \Needspace{35mm}\BankTaskContent{#1}{#2}{#3}%
  \else
    \Needspace{\dimexpr\ht\BankTaskBox+\dp\BankTaskBox\relax}%
    \noindent\usebox{\BankTaskBox}\par
  \fi
  \medskip
}
\begin{document}
\BookSubsection{6.11.2}{Деление ядер}


% Source 6E1C06: https://ege.fipi.ru/bank/questions.php?search=1&proj=BA1F39653304A5B041B656915DC36B38&qid=6E1C06
\BankTask{6E1C06}{
В результате цепной реакции деления урана \({}_{0}^{1}\mathrm{n}+{}_{92}^{235}\mathrm{U}\to X+{}_{36}^{94}\mathrm{Kr}+3{}_{0}^{1}\mathrm{n}\) образуется ядро химического элемента \({}_{Z}^{A}X\). Каковы заряд \(Z\) образовавшегося ядра (в единицах элементарного заряда) и его массовое число \(A\)?
\begin{center}\begin{tabular}{|c|c|}\hline
Заряд ядра $Z$ & Массовое число ядра $A$ \\ \hline
\rule{0pt}{1.5em} & \\ \hline
\end{tabular}\end{center}
}{
\smallskip\noindent\textbf{Решение от GPT}\par
По закону сохранения заряда \(Z=92-36=56\). По закону сохранения массового числа \(A=1+235-94-3=139\). Ответ: \(Z=56\), \(A=139\).
}

\end{document}
